\documentclass[aspectratio=169]{beamer}
\input{header}
\coursecode{JKSSB}

\title{Cybersecurity, IT Act \& Security Controls}
\subtitle{Lesson 56 --- Law, Threats and Controls}
\author{JKSSB Computer Awareness}
\date{\today}

\begin{document}
\frame{\titlepage}

\begin{frame}{Why Security Matters: The CIA Triad}
  \begin{itemize}
    \item A computer holds data worth protecting: marks, accounts,
      personal records, and official files.
    \item \key{Security} means protecting data and systems against theft,
      damage, and misuse.
    \item Three goals, remembered as \key{CIA}: \key{confidentiality}
      (only authorised persons read), \key{integrity} (data is not
      altered), \key{availability} (systems stay usable).
    \item Every control in this lesson serves one or more of these goals.
  \end{itemize}
  \begin{block}{The one idea}
    Threats attack the CIA triad; controls defend it. Match each control
    to what it does.
  \end{block}
\end{frame}

\begin{frame}{The IT Act, 2000}
  \begin{itemize}
    \item The \key{Information Technology Act, 2000} is India's law dealing
      with electronic records, electronic governance, and cybercrime.
    \item It gives \key{legal validity} to electronic records and electronic
      signatures, so online transactions and e-governance have a legal basis.
    \item It defines \key{cyber contraventions and offences}, such as
      unauthorised access, damage to computer systems, and fraud committed
      through a computer.
    \item Year and scope are the exam facts: the Act is of \key{2000} and
      covers \key{electronic records, e-governance, and cybercrime}.
  \end{itemize}
  \begin{alertblock}{Trap alert}
    The IT Act is of \key{2000}, not 2008 as a new Act --- later changes
    were amendments to the 2000 Act. Do not invent section numbers.
  \end{alertblock}
\end{frame}

\begin{frame}{The Malware Umbrella}
  \begin{itemize}
    \item \key{Malware} (malicious software) is the umbrella term for any
      program written to harm a computer, steal data, or disrupt work.
    \item The four named members: \key{virus}, \key{worm},
      \key{Trojan horse}, and \key{ransomware}.
    \item Related nuisance programs: \key{spyware} (secretly watches and
      reports activity) and \key{adware} (shows unwanted advertisements).
    \item Exam wording: malware is the \key{general} term; virus, worm,
      Trojan, and ransomware are \key{specific} types under it.
  \end{itemize}
  \begin{block}{Keep the pair straight}
    Malware is the umbrella; each named threat is one rib of the umbrella.
  \end{block}
\end{frame}

\begin{frame}{Virus: Needs a Host and Human Action}
  \begin{itemize}
    \item A \key{virus} attaches itself to a host file or program and
      spreads when that file is run or shared.
    \item It \key{needs human action} --- opening an infected attachment or
      running an infected file --- to spread.
    \item It can corrupt, delete, or alter files once it runs.
    \item Office-desktop picture: an infected document copy passed on a
      pen drive infects the next machine that opens it.
  \end{itemize}
  \begin{alertblock}{Exam contrast}
    Virus needs a \key{host file plus human action}; worm does not.
  \end{alertblock}
\end{frame}

\begin{frame}{Worm: Self-Spreading Across Networks}
  \begin{itemize}
    \item A \key{worm} is standalone malware that \key{copies itself} and
      spreads across a network \key{without} a host file.
    \item It needs \key{no human action} to move from machine to machine;
      it exploits the network itself.
    \item Worms consume \key{bandwidth and system resources}, slowing or
      crashing networks.
    \item \muted{One infected machine can flood the office network while
      nobody opens anything.}
  \end{itemize}
  \begin{block}{Virus vs worm}
    Virus = host file + human action. Worm = self-copying, no host, no
    click needed.
  \end{block}
\end{frame}

\begin{frame}{Trojan Horse: Looks Useful, Hides Harm}
  \begin{itemize}
    \item A \key{Trojan horse} pretends to be useful software but hides a
      harmful purpose inside.
    \item Unlike a virus or worm, it \key{does not replicate} by itself.
    \item The harm: it opens a \key{backdoor}, letting an attacker steal
      passwords, spy on activity, or control the machine.
    \item Office-desktop picture: a ``free'' utility that installs fine ---
      and quietly hands the attacker the keyboard.
  \end{itemize}
  \begin{alertblock}{Trap alert}
    Trojan does \key{not} self-replicate. Replication belongs to virus
    and worm.
  \end{alertblock}
\end{frame}

\begin{frame}{Ransomware: Locks, Then Demands Payment}
  \begin{itemize}
    \item \key{Ransomware} \key{encrypts} the victim's files or locks the
      screen, then demands payment (\key{ransom}) for release.
    \item Payment does \key{not} guarantee recovery; the defence is a clean
      \key{backup} kept separate from the machine.
    \item It commonly arrives through phishing mail and infected downloads.
    \item \muted{Exam one-liner: ransomware = locks data + demands money.}
  \end{itemize}
  \begin{exampleblock}{Worked contrast}
    Virus corrupts files; worm floods the network; Trojan spies through a
    backdoor; ransomware locks files and demands payment.
  \end{exampleblock}
\end{frame}

\begin{frame}{Antivirus and Quarantine}
  \begin{itemize}
    \item \key{Antivirus} software detects, blocks, and removes malware,
      using known malware signatures and suspicious-behaviour checks.
    \item It must be \key{updated regularly}, since new malware appears
      every day.
    \item \key{Quarantine} means isolating a suspicious file so it cannot
      run or spread while awaiting a decision.
    \item Quarantine is \key{isolation}, not deletion and not a backup.
  \end{itemize}
  \begin{block}{Keep the pair straight}
    Delete removes; quarantine isolates. An isolated file can still be
    inspected or restored.
  \end{block}
\end{frame}

\begin{frame}{Phishing and Anti-Phishing}
  \begin{itemize}
    \item \key{Phishing} is a fake message (mail, SMS, or site) that
      impersonates a trusted sender to steal passwords or money.
    \item Warning signs: urgent threats, unexpected attachments, links that
      do not match the claimed sender, requests for OTPs or passwords.
    \item \key{Anti-phishing} defences: verify the sender, hover over links
      before clicking, never share OTPs, and report the message.
    \item A genuine bank or office \key{never} asks for a password or OTP
      by mail or message.
  \end{itemize}
  \begin{alertblock}{Trap alert}
    Phishing steals \key{credentials by deception}; it is not a virus and
    needs no malware to succeed.
  \end{alertblock}
\end{frame}

\begin{frame}{Encryption vs Digital Signature}
  \begin{itemize}
    \item \key{Encryption} converts readable data (plaintext) into coded
      form (ciphertext) so only authorised persons can read it ---
      it gives \key{confidentiality}.
    \item A \key{digital signature} proves who sent a message and that it
      was not altered --- it gives \key{authenticity and integrity}.
    \item A signature does \key{not} hide the message; encryption does.
    \item \muted{This repeats the Lesson-16 email rule (TRAP-15): signature
      $\ne$ encryption.}
  \end{itemize}
  \begin{block}{One line each}
    Encryption hides content. Signature proves sender and intactness.
  \end{block}
\end{frame}

\begin{frame}{SSL and TLS: Secure Transmission}
  \begin{itemize}
    \item \key{SSL} stands for Secure Sockets Layer; \key{TLS} stands for
      Transport Layer Security.
    \item TLS is the \key{successor} of SSL; the pair provides a secure,
      encrypted channel over a network.
    \item They protect data \key{in transit}: mail sending, web login, and
      online forms (the padlock in the browser address bar).
    \item Exam anchor (WO 2026 Q76): the secure email transmission protocol
      is \key{SSL/TLS}.
  \end{itemize}
  \begin{alertblock}{Trap alert}
    SSL/TLS protect data in transit. They do not clean an infected file
    and do not replace antivirus.
  \end{alertblock}
\end{frame}

\begin{frame}{Firewall: The Network Barrier}
  \begin{itemize}
    \item A \key{firewall} is a barrier that filters incoming and outgoing
      network traffic against a set of rules.
    \item It \key{blocks unauthorised access} while letting permitted
      traffic pass.
    \item It can be \key{hardware}, \key{software}, or both.
    \item \muted{Picture it as the gatekeeper at the office network door,
      checking every packet against the list.}
  \end{itemize}
  \begin{exampleblock}{Worked contrast}
    Firewall filters \key{traffic} at the boundary; antivirus cleans
    \key{malware} inside the machine. They complement, not replace, each
    other.
  \end{exampleblock}
\end{frame}

\begin{frame}{Patching and MFA: Two More Controls}
  \begin{itemize}
    \item \key{Patching} means installing vendor-supplied software updates
      that fix security flaws; an unpatched machine stays open to known
      attacks.
    \item \key{MFA} stands for Multi-Factor Authentication: proving identity
      with \key{two or more} checks, e.g.\ password plus OTP or fingerprint.
    \item MFA protects \key{logins}; patching protects \key{software flaws}.
    \item Strong, distinct \key{passwords} remain the first factor ---
      MFA adds factors, it does not replace the password.
  \end{itemize}
  \begin{block}{Which control does what}
    Patching closes holes in software. MFA closes the door against stolen
    passwords.
  \end{block}
\end{frame}

\begin{frame}{Which Control Does What: Matching Map}
  \begin{center}
    \begin{tabular}{ll}
      \toprule
      \textbf{Control} & \textbf{Job} \\
      \midrule
      Firewall & Filters network traffic; blocks unauthorised access \\
      Antivirus & Detects and removes malware on the machine \\
      Quarantine & Isolates a suspicious file so it cannot run \\
      Encryption (SSL/TLS) & Protects data in transit from reading \\
      Digital signature & Proves sender and intactness \\
      Patching & Fixes security flaws in software \\
      MFA & Adds extra login checks beyond the password \\
      Backup & Recovers data after ransomware or loss \\
      \bottomrule
    \end{tabular}
  \end{center}
  \vspace{0.25cm}
  \muted{Read the stem's verb --- filters, cleans, isolates, hides,
  proves, fixes, or recovers --- then pick the matching row.}
\end{frame}

\begin{frame}{Trap Alert: Security Facts to Keep Straight}
  \begin{itemize}
    \item IT Act = \key{2000}; covers electronic records, e-governance,
      and cybercrime.
    \item Malware is the \key{umbrella}; virus, worm, Trojan, and
      ransomware are types under it.
    \item Virus needs host + action; worm self-spreads; Trojan does
      \key{not} replicate; ransomware locks and demands payment.
    \item Digital signature $\ne$ encryption: signature proves,
      encryption hides. (TRAP-15)
    \item SSL/TLS = secure \key{transmission}; firewall = traffic
      \key{barrier}; quarantine = \key{isolation}; MFA = \key{extra login
      checks}.
  \end{itemize}
\end{frame}

\begin{frame}{Lesson 56: Recall}
  \begin{itemize}
    \item The Information Technology Act, 2000 gives legal validity to
      electronic records and signatures and defines cyber contraventions
      and offences.
    \item Malware is the umbrella term; virus (host + action), worm
      (self-spreading), Trojan (false-useful backdoor, no replication),
      and ransomware (locks + demands payment) sit under it.
    \item Antivirus detects and removes malware; quarantine isolates a
      suspicious file; patching fixes software flaws.
    \item Phishing steals credentials by deception; verify senders, never
      share OTPs.
    \item Encryption gives confidentiality; digital signature gives
      authenticity and integrity; SSL/TLS secure data in transit;
      firewall filters traffic; MFA adds login factors.
  \end{itemize}
\end{frame}
\end{document}
